% Pro A. Licinio Archia Poeta (pro-archia) --- English translation
% Translated by Claude (Anthropic) from the Latin (Perseus).
% Style guide: STYLE.md.

\ciceroSection{1} If there is in me, judges, any talent --- and I feel how slight it is --- or any practice in speaking, in which I do not deny that I have been moderately exercised, or any reasoned method in this art that springs from the study and discipline of the best learning, from which at no period of my life I confess I have ever held aloof, then for all these things this Aulus Licinius is almost by his own right entitled to claim from me the fruit, and that among the first. For as far back as my mind can look upon the stretch of time gone by, and recollect the furthest memory of my boyhood, I see that, going back to that very beginning, this man stood out for me as the foremost in undertaking and entering upon the path of these studies. So if this voice of mine, formed by his exhortation and instruction, has at times been of some help to others, surely to him from whom I received what we have been able to use to assist others and save some, we ought, so far as it lies in us, to bring help and safety in turn.

\ciceroSection{2} And lest anyone wonder, perhaps, that I say this --- because the sort of talent in him is of another kind, neither this method nor this discipline of speech --- let me say that even I have never been wholly given up to this one pursuit. For all the arts that pertain to humane culture have a certain bond in common, and are held together as if by a kinship among themselves.

\ciceroSection{3} But that it may not seem strange to any of you that, in a regular inquiry, in a public trial, when the matter is being treated before a praetor of the Roman people, that most upright of men, and before the strictest of judges, with so great a gathering and crowd of men, I should use this kind of speech which differs not only from the custom of the courts but even from forensic speech, I beg of you that in this case you grant me this indulgence, suited to this defendant and (as I hope) not unwelcome to you: that, while I speak for a poet of the highest rank and a man of the deepest learning, in this gathering of the most cultivated men, with this humanity of yours, with this praetor presiding, you suffer me to speak somewhat more freely about the studies of humane culture and letters, and in the case of a person of this kind --- one who, on account of his quiet and his studies, has been very little tried in the courts and in the perils of trial --- to use a kind of speech that is almost new and unaccustomed.

\ciceroSection{4} And if I see that this is being granted and conceded to me by you, I shall surely make you think that this Aulus Licinius, since he is a citizen, must not be set apart from the number of citizens, but even, were he not one, must have been received into it. For as soon as Archias outgrew boyhood and turned himself, from those arts in which boys are wont to be shaped to humane culture, to the pursuit of writing, first at Antioch --- for there he was born of noble stock --- once a celebrated city, rich and abounding in men of the deepest learning and in the most liberal pursuits, he quickly began to outstrip all in glory of talent. Afterwards in the rest of the parts of Asia and through all Greece his arrivals were so celebrated that the expectation of the man surpassed the fame of his talent, and the wonder of his very arrival surpassed the expectation.

\ciceroSection{5} Italy then was full of Greek arts and disciplines, and these studies in Latium too were more vehemently cultivated then than they are now in the same towns; and here at Rome, on account of the tranquillity of the commonwealth, they were not neglected. So both the Tarentines and the men of Locri and the Rhegines and the Neapolitans bestowed on him citizenship and other gifts, and all who could give any judgment on talents reckoned him worthy of acquaintance and of hospitality. With his fame so widely celebrated, when he was already known even to those who had never met him, he came to Rome in the consulship of Marius and Catulus. He found, first, consuls of whom the one could supply him with the greatest matter for writing, the other could bring not only deeds done but enthusiasm and an ear for letters. At once the Luculli, while Archias was still wearing the praetexta, took him into their house. And this was a tribute given not only to the brightness of his talent and his learning but to his nature and virtue: that the house which had favoured this man's youth was the same that became most familiar to his old age.

\ciceroSection{6} He was at those times welcome to that famous Quintus Metellus Numidicus and to his son Pius; he was listened to by Marcus Aemilius; he lived with Quintus Catulus, both father and son; he was cultivated by Lucius Crassus; and since he held the Luculli, Drusus, the Octavii, Cato, and the whole house of the Hortensii bound to him by daily intercourse, he was treated with the highest honour. For not only those zealous to learn and listen cultivated him, but also any who happened to pretend so. Meanwhile, after a long enough interval, when he had set out for Sicily with Marcus Lucullus, and when from that province he was leaving with the same Lucullus, he came to Heraclea. Since that city was on the most equal terms by treaty, he wished to be enrolled in that citizenship, and what was reckoned his by his own right he obtained from the Heracleans by the authority and favour of Lucullus.

\ciceroSection{7} Citizenship was given by the law of Silvanus and Carbo: ``If any have been enrolled in confederate cities, if at the time when the law was being carried they had a domicile in Italy, and if within sixty days they had registered before the praetor.'' Since this man had had his domicile at Rome for many years already, he registered before the praetor

\ciceroSection{8} Quintus Metellus, his closest friend. If we are speaking of nothing else but the citizenship and the law, I say no more: the case has been pleaded. For what of these points, Grattius, can be shaken? Will you deny that he was then enrolled at Heraclea? There is present a man of the highest authority, scrupulousness, and good faith, Marcus Lucullus; who says that he does not hold an opinion but knows, did not hear but saw, was not present but acted. Present too are the legates of Heraclea, men of the highest standing; for the sake of this trial they have come with their instructions and with the public testimony of their city, and say that this man was enrolled at Heraclea. Here you demand the public records of Heraclea --- which, as we all know, perished when the records office was burned in the Italian war? It is ridiculous to have nothing to say to what we have, and to ask after what we cannot have; to keep silent about the memory of men, and to demand the memory of letters; and, when you have the scruple of a man of the highest standing and the oath and good faith of a town of the most absolute integrity, to reject the things which can in no way be falsified and to demand records which, as you yourself say, are wont to be tampered with.

\ciceroSection{9} Or did he not have a domicile at Rome --- the man who, so many years before the citizenship was given, set the seat of all his affairs and fortunes at Rome? Or did he not register? On the contrary: he registered in those very records which alone, of that registration and that college of praetors, hold the authority of public records. For when the records of Appius were said to have been kept too negligently, and when the levity of Gabinius (so long as he was unbroken) and his disaster after his condemnation had unsealed all credit in his records, Metellus, the most scrupulous and most modest of men, was of such diligence that he came to the praetor Lucius Lentulus and to the judges and said that he was disturbed by the erasure of a single name. In these records, then, you see no erasure on the name of Aulus Licinius.

\ciceroSection{10} Since this is so, what is there for you to doubt about his citizenship --- especially when in other states too he was enrolled? For in fact, when the Greeks would impart citizenship gratuitously to many men of middling stature, endowed with no art or with some humble one, are we to suppose that the Rhegines or Locrians or Neapolitans or Tarentines, what they used to lavish on stage performers, did not wish to bestow on this man, endowed with the highest glory of talent? What? When others have crept somehow into the records of those towns not only after the citizenship had been given but even after the Lex Papia, shall this man --- who does not even use those rolls in which he is enrolled, because he has always wished to be a Heraclean --- be cast out?

\ciceroSection{11} You ask after our census-rolls. Of course --- for it is no secret that at the most recent census he was with the most distinguished commander Lucius Lucullus on his army; at the previous one with the same Lucullus as quaestor in Asia; at the first, when Julius and Crassus were censors, no part of the people was assessed. But, since the census does not confirm the right of citizenship and only shows that the man who was assessed had even then conducted himself as a citizen, in those very times the man you accuse --- not even by his own judgment --- of standing in the right of Roman citizens both made his will many times under our laws, and entered upon the inheritances of Roman citizens, and was put up to the treasury for benefits by Lucius Lucullus the proconsul. Look for arguments, if you have any: this man will never be refuted either by his own or by his friends' judgment.

\ciceroSection{12} You will ask of us, Grattius, why I take such great delight in this man. Because he supplies us with that wherewith both the mind may be refreshed from this forensic clamour, and the ears, weary with brawling, may rest. Or do you suppose either that I could find what to say day after day, in such great variety of cases, did I not till my mind by learning, or that my mind could bear so great a strain, did we not relax it by that same learning? I confess: I am given over to these studies. Let the rest be ashamed --- those who have so buried themselves in letters that they can bring nothing from them either to the common fruit or out into sight and the light. But why should I be ashamed, who for so many years live in such a way, judges, that no man's crisis or interest has ever drawn me from my leisure, or any pleasure called me away, or, finally, any sleep delayed me?

\ciceroSection{13} Who, then, would rebuke me, or who would justly be angry with me, if all the time which is conceded to others for going about their own business, for celebrating the festal days of the games, for other pleasures and for the very rest of mind and body, all that others spend on protracted dinner-parties, finally on the gaming-board, on ball-play --- as much of that as that, I take back to myself for the renewing of these studies? And this should be granted to me the more, that out of these studies grows also this oratorical faculty, which (such as it is in me) has never failed my friends in their dangers. And if to anyone this seems a slight matter, I at least feel from what fount I draw those things which are highest.

\ciceroSection{14} For unless from my youth I had persuaded myself, by the precepts of many and by much reading, that nothing in life is much to be sought after but praise and honourable conduct, and that in pursuing this all bodily torments, all dangers of death and exile are to be reckoned light, never should I have thrown myself for your safety into so many and so great battles, and into these daily attacks of ruined men. But all the books are full of this, full are the voices of the wise, full of examples is antiquity --- and all this would lie in darkness, did not the light of letters draw near. How many images of the bravest of men our writers, both Greek and Latin, have left behind for us, fashioned not only for our looking upon, but for our imitation! These I, in administering the commonwealth, always set before me, and I formed my mind and spirit by the very contemplation of men of excellence.

\ciceroSection{15} Someone will ask: ``What? Were those very greatest of men, whose virtues have been handed down in letters, learned in that very discipline you so extol with praises?'' It is hard to say so of them all, but yet what I should reply is sure. I confess that there have been many men of excellent spirit and virtue without learning, who, by some almost divine cast of nature itself, of their own accord, stood out as both moderate and weighty men. I add this also: more often nature without learning has availed for praise and virtue than learning without nature. And this same point I press: when to a nature outstanding and bright there is added some method and shaping of learning, then that I-know-not-what bright and singular thing is wont to come into being.

\ciceroSection{16} Of this number was the man whom our fathers saw, that godlike man Africanus; of this Gaius Laelius, Lucius Furius, the most moderate and continent of men; of this the bravest of men and, in those times, the most learned, that elder Marcus Cato. Surely if these had not been helped by letters in any way for the grasping and cultivating of virtue, they would never have given themselves to their study. But even if no such great fruit were shown, and only delight were sought from these studies, yet, in my view, you would judge this relaxation of the mind the most humane and the most liberal of all. For other relaxations belong neither to all times nor to all ages nor to all places: but these studies sharpen youth, delight age; they adorn prosperity, in adversity afford a refuge and consolation; they delight us at home, do not hinder us abroad; they pass the night with us, they travel, they go to the country.

\ciceroSection{17} And even if we ourselves could neither touch these things nor taste them with our own senses, yet we ought to admire them when we saw them in others. Who of us was of so rude and hard a spirit that he was not lately moved by the death of Roscius? When he had died an old man, yet because of his outstanding art and grace he seemed not at all to have had to die. So that man, by the movement of his body, had drawn so great a love to himself from us all: shall we neglect the incredible movements of mind and the swiftness of talents?

\ciceroSection{18} How often have I seen this Archias, judges --- I shall use your indulgence, since you attend to me so diligently in this new kind of speech --- how often have I seen this man, when he had not put pen to paper, deliver on the spot a great quantity of excellent verses on the very matters then being treated! How often, when called upon, deliver the same matter again with words and thoughts changed! And what he had carefully and thoughtfully written, I have seen approved as to attain the praise of the ancient writers. Shall I not love this man, not admire him, not think him in every way to be defended? And we have learned this from the highest and most learned men: that the studies of other things stand on learning and precept and art, but the poet is strong of nature itself, and is roused by the powers of mind, and is breathed into, as it were, by some divine spirit. Wherefore by his own right our Ennius calls poets ``holy,'' because they seem to have been commended to us by some gift, as it were, and grant of the gods.

\ciceroSection{19} Let it then, judges, be holy with you, the most cultured of men, this name of poet, which no barbarism has ever violated. The rocks and the wildernesses answer to the voice; savage beasts are often softened by song and stand still. Shall we, brought up on the best things, not be moved by the voice of poets? The Colophonians say that Homer is their citizen, the Chians claim him as theirs, the Salaminians ask him back, but the Smyrnaeans assert that he is theirs --- and for that reason have even consecrated a shrine to him in their town; and very many others besides fight and contend among themselves. They, then, claim a man not of their own (because he was a poet) even after his death; shall we reject this man, alive, who is ours both by his own will and by our laws --- especially since once Archias bestowed his whole pursuit and his whole talent on celebrating the glory and praise of the Roman people? For as a young man he treated of the Cimbric war, and was welcome to that very Gaius Marius, who seemed harder for these studies.

\ciceroSection{20} For there is no one so averse from the Muses that he would not gladly let an eternal proclamation of his labours be entrusted to verses. Themistocles, that greatest man at Athens, is said to have answered, when he was asked what music or whose voice he most gladly heard: ``the voice of him by whom my own virtue may best be declared.'' So that same Marius likewise singularly loved Lucius Plotius, by whose talent he thought what he had done could be celebrated.

\ciceroSection{21} Indeed the Mithridatic war, great and difficult and engaged in much variety by land and sea, has been wholly set forth by him; which books not only illuminate Lucius Lucullus, the bravest and most distinguished of men, but the very name of the Roman people. For the Roman people opened up Pontus, with Lucullus in command --- a Pontus once walled in by royal wealth and by nature itself and by the lie of the land. The army of the Roman people, with the same leader, with no great force routed innumerable hordes of Armenians. It is the praise of the Roman people that the most friendly city of the Cyzicenes was, by his counsel, snatched from every royal assault and from the very mouth and jaws of the whole war, and saved. Ours always will be borne and celebrated, while Lucullus fought, the incredible naval battle at Tenedos, where, with the leaders killed, the enemy's fleet was sunk; ours are the trophies, ours the monuments, ours the triumphs. By whose talents these things are set out, by them the fame of the Roman people is celebrated.

\ciceroSection{22} Dear to the elder Africanus was our Ennius, and so even in the tomb of the Scipios he is thought to have been set up in marble. But by those praises certainly not only is he himself who is praised adorned, but the name of the Roman people too. To the heavens is borne by them his great-grandfather Cato; great honour is added to the deeds of the Roman people. All those greatest men, finally, the Marcelli, the Fulvii, are decorated not without the common praise of all of us. Therefore that man who had done these things, the man of Rudiae, our ancestors received into citizenship: shall we drive out from our citizenship this Heraclean, sought after by many cities and in this established by the laws?

\ciceroSection{23} For if anyone thinks that less fruit of glory is reaped from Greek verses than from Latin, he is vehemently in error, since Greek is read among nearly all peoples, while Latin is shut within its own bounds, indeed narrow ones. Wherefore, since the things we have done are bounded by the regions of the world, we ought to wish that wherever the weapons of our men have reached, our glory and fame should reach the same place; both because this is itself a great thing for those very peoples about whose deeds is written, and because it is the strongest spur, for those who fight at risk of life for glory's sake, both to dangers and to labours.

\ciceroSection{24} How many writers of his deeds is that great Alexander said to have had with him! And yet, when he had stood at Sigeum at the tomb of Achilles, ``O fortunate young man,'' he said, ``who hast found in Homer the herald of thy virtue!'' And rightly. For if that Iliad had not come to be, the same mound that had covered his body would have buried his name too. What of him? Did our own Magnus, who has matched virtue with fortune, not give Theophanes of Mitylene, the writer of his deeds, citizenship at a meeting of soldiers? And those brave men of ours, but country-folk and soldiers as they were, moved by some sweetness of glory --- as if they themselves were sharers in the same praise --- approved this with a great shout?

\ciceroSection{25} So, I take it, if Archias were not a Roman citizen by the laws, he could not have brought it about that he should be given citizenship by some commander. Sulla, when he gave it to Spaniards and Gauls, would have rejected this man's request, I take it. Whom we saw, when a bad poet from the people had thrust a little book at him --- because he had made an epigram about him only in alternating, slightly longer verses --- at once order, out of those very things which he was then selling, that a reward be given him, but on this condition: that he write nothing afterwards. Should the man who reckoned the busy effort of a bad poet worth some reward not have sought after this man's talent and virtue and copiousness in writing?

\ciceroSection{26} What? Would he not have got it from Quintus Metellus Pius, his closest friend, who gave citizenship to many, either through himself or through the Luculli? --- who, indeed, was so eager that his deeds should be written about that he gave his ear even to poets born at Corduba, who had something fat and foreign about their sound. For there is no concealing what cannot be hidden, but it must be carried before us: we are all drawn by the love of praise, and the best men most of all are led by glory. The very philosophers themselves, even in those little books which they write about despising glory, inscribe their own names; in that very thing in which they despise being talked of and noised about, they wish to be talked of and to be named.

\ciceroSection{27} Decimus Brutus indeed, that greatest of men and commander, adorned the entrances of the temples and monuments of his family with the songs of his closest friend Accius. And that Fulvius who waged war with the Aetolians, with Ennius as his companion, did not hesitate to consecrate to the Muses the spoils of Mars. Wherefore, in a city in which commanders almost in arms have cultivated the name of poets and the shrines of the Muses, in that city judges in the toga ought not to shrink from honour to the Muses and from the safety of poets.

\ciceroSection{28} And, that you may do this the more gladly, I shall now declare myself to you, judges, and confess to you a certain love of glory in me, perhaps too keen, but yet honourable. For those things which we, in our consulship, did along with you for the safety of this city and empire, for the lives of citizens, and for the whole commonwealth, this man has touched in verses and begun. Hearing them, because the matter seemed to me great and welcome, I encouraged him to bring it to completion. For virtue desires no other wages of its labours and dangers but this of praise and glory; with which taken away, judges, what is there for which in this so brief and narrow course of life we should exercise ourselves in such great labours?

\ciceroSection{29} Surely, if the mind foresaw nothing for the future, and if by the same regions in which the span of life is bounded it also bounded all its thoughts, neither would it break itself with so great labours nor be vexed by so many cares and watchings, nor so often fight for life itself. Now there sits in every best man a certain virtue which night and day stirs the mind with the goads of glory, and reminds us that the commemoration of our name must not be let go with the time of our life, but must be made equal with all posterity.

\ciceroSection{30} Or are we all to seem so small of mind --- we who are engaged in the commonwealth and in these dangers and labours of life --- that, when up to our last span we have drawn no tranquil and quiet breath, we are to think that all things will die together with us? But have many great men diligently left behind statues and likenesses --- not images of the mind, but of the body --- and ought we not far more to wish to leave behind a likeness of our counsels and our virtues, set forth and polished by the highest talents? For my part I thought, in the very doing, that everything I was doing was being scattered and disseminated by me into the everlasting memory of the world. Whether after death this will be wanting to my sense, or whether (as the wisest of men have thought) it will pertain to some part of my mind, now at any rate I take pleasure in some thought and hope of it.

\ciceroSection{31} Wherefore, judges, save a man of such modesty as you see approved by his friends together with their dignity and their long acquaintance, of so great talent as it befits to be reckoned --- talent which you see has been sought after by the judgments of the highest men, of a case such as is approved by the benefit of a law, by the authority of a town, by the testimony of Lucullus, by the records of Metellus. Since this is so, we beg of you, judges, if there ought to be in such great talents any commendation not only human but even divine, that this man --- who has always adorned you, your commanders, the deeds of the Roman people, who professes that of these recent dangers, mine and yours and our country's, he will give an eternal testimony of praise, and who is of that number which has always been held holy and so called among all men --- you take into your faith in such a way that he may seem rather to have been relieved by your humanity than wounded by your harshness.

\ciceroSection{32} What, after my custom, I have said briefly and simply on the case, judges, I trust has been approved by all. What I have said, foreign to the Forum and the custom of the courts, both about the man's talent and in general about the pursuit itself, this, judges, I hope has been taken in good part by you; and by the man who presides over the trial, I am sure of it.



